\subsection{有界输入批次与两遍构建}

一次旋转全部 $N$ 个原始向量需要两个 FP32 设备数组，共 $8Nd$ 字节。构建器改用 pinned 主机内存和设备原始／旋转缓冲区分批暂存。若有 $n_b$ 个流水缓冲区、配置 add 行数 $a_{\rm add}$、每数组字节预算
$S_{\rm stage}$，源码暂存行数为
\begin{equation}
 a_{\rm in}=\max\!\left\{1,\min\!\left(a_{\rm add},
 \left\lfloor\frac{S_{\rm stage}}{4d}\right\rfloor\right)\right\},
 \qquad S_{\rm float}=8n_ba_{\rm in}d.
 \label{eq:4-staging-rows}
\end{equation}
默认三个缓冲区、128\,MiB 预算，使输入／旋转设备浮点数组至多占 768\,MiB。
分配乘积和编码 scratch 是独立分配。

一遍路径先编码输入顺序的主码／残差数组，再 gather 成列表顺序，临时与最终压缩副本共存。容量规则预测内存不足时，两遍构建先分配列表、排序 ID、构建偏移，再次读取输入，旋转、编码并直接 scatter 到最终位置。它用额外一次输入扫描和旋转换取较少的同时存活压缩副本。两遍路径的默认分配可在原始输入与反变换中心之间比较；正交性给出相应的实数算术恒等式。

\subsection{阶段占用与进程峰值}

源码的一遍路径分支估计为
\begin{equation}
 H_{\rm single}=N(2b_r+3M+20)+2^{30}\quad\text{bytes}.
 \label{eq:4-peak-heuristic}
\end{equation}
其中依赖 $N$ 的项，是 gather／transpose 时同时存活的十个数组的精确载荷：输入顺序的主码、残差和修正项；分配结果及排序顺序；最终 ID、残差和修正项；vector-major 与 subspace-major 主码副本。共计 $3NM+2Nb_r+20N$ 字节。额外 1\,GiB 是容量启发式的预留，不是其他全部分配的精确核算。

两遍构建在最终 repack 时，同时保留 ID、残差、修正项和新旧主码数组：
\begin{equation}
 S_{\rm repack,2}=N(b_r+2M+8).
 \label{eq:4-two-pass-repack}
\end{equation}
这一表达式对该阶段的这些数组精确成立。更早阶段还保留流水线 scratch，scatter 时另有临时逆置换。因此，repack 表达式本身不是两遍构建的整体峰值。取整个生命周期最大值时，还需计入模型、排序／分配器 scratch、运行时上下文及之后的工作区分配。

\begin{table}[tbp]
\centering
\caption[各构建阶段的行数组显存模型]{行数组的解析占用（GiB），$d=1{,}024$、$M=128$、$B_r=8$。BGE 使用一遍，Stella 使用两遍；Stella 一遍列为反事实容量分析。这些阶段值不含流水线、模型、偏移和运行时成本。}
\label{tab:4-phase-storage}
\small
\begin{tabular}{@{}lrrr@{}}
\toprule
数据集 & 最终行载荷 & 一遍 gather & 两遍 repack \\
\midrule
BGE & 10.90 & 23.05 & 12.11 \\
Stella & 19.20 & 40.59 & 21.32 \\
\bottomrule
\end{tabular}
\end{table}

\Cref{tab:4-phase-storage}量化了压缩副本的节省。Stella 冷日志选择两遍，报告\cref{eq:4-peak-heuristic}为 41.6\,GiB，尾部可用为 30.6\,GiB。另行采样的常驻进程总量包括上下文及查询工作区，BGE 为 11.93\,GiB，Stella 为 20.44\,GiB。容量模型见\cref{sec:4-4}。
补证以20\,ms间隔采样，BGE在add开头观测26.08\,GiB，Stella在repack阶段
观测22.58\,GiB。这是含上下文的设备采样最大值，可能漏掉更短暂分配。

分配乘积容量使用独立配置行数 $a_{\rm assign}$，由
\texttt{JHQ\_ASSIGN\_BATCH} 或默认值选择：列表数至少 8,192 时为 32,768，
否则为 8,192。
\begin{equation}
 S_{\rm dots}=4n_b n_{\rm list}a_{\rm assign}\quad\text{bytes}.
 \label{eq:4-assignment-scratch}
\end{equation}
分配容量不必随 $a_{\rm in}$ 或某次调用的实际处理行数缩小。
\Cref{tab:b-scratch}核对四种默认配置：1,536／3,072 维虽暂存行数不同，
分配乘积均占 3\,GiB。BGE／Stella 在 add 入口附近仍分配三个各 4\,GiB 数组，
共 12\,GiB。
23.05\,GiB的gather载荷属于后续阶段，不能直接再加12\,GiB作为真实峰值。
峰值模型应对各阶段实际存活数组（含流水scratch）取最大值。
Stella两遍repack的较小载荷加模型／上下文，解释了其数据库更大却观测峰值更低。

\subsection{构建测量}

\Cref{sec:7-7}通过重复空缓存构建、较宽启动区间和采样设备峰值评价上述阶段模型，
并报告基线构建与大数据集搜索。计时边界沿用\cref{sec:3-3}；
上述构建算法和分配生命周期决定哪些成本可以重叠。
